\section{Analysis and discussion}
\label{sec:analysis}

The results are consistent with the Azure tokenizer allocating vocabulary to
domain-recurring strings and documentation patterns, so that resource-provider
names, command flags, paths, and configuration keys are encoded in fewer pieces.
The mechanism appears in fragmentation and vocabulary-utilization measurements in
the repository: on the principal holdout the Azure tokenizer encodes
resource-provider identifiers, URLs, command flags, and camel/snake identifiers in
fewer tokens than the General tokenizer (for example, roughly half as many tokens
per resource-provider identifier in the repository's fragmentation table), and it
uses a larger fraction of its vocabulary on Azure text.

Several observations bound the interpretation. First, the effect \emph{persists on
prose} (the largest content stratum, $n=900$, still shows \reduc{+36.75}), which
weakens the alternative that gains come only from code fences or configuration
syntax. Second, the advantage \emph{transfers} to non-Microsoft Kubernetes
documentation (\reduc{+31.18}), suggesting overlap in general
technical-documentation style rather than only Azure-specific terminology. Third,
the reduction is \emph{stable across} 8K/16K/32K, suggesting corpus composition
mattered more than the tested vocabulary range. Fourth, the Mixed tokenizer is the
central practical result: it recovers most of the in-domain gain with a small
out-of-domain penalty, which matters when a single tokenizer must serve mixed
traffic.

We are explicit about a design limit: comparing Azure documentation with general
English changes \emph{more than subject domain}. It also changes technical writing
style, markup distribution, identifier density, and document structure. The
transfer and prose-only analyses narrow, but do not fully isolate, Azure-specific
vocabulary from general technical-documentation specialization
(Section~\ref{sec:limits}). Finally, faster compression does not imply faster CPU
tokenization---the controlled tokenizers compress more yet run slower on CPU than
the production baselines---and none of these measurements support any downstream
model claim.

\subsection*{What this study is and is not}

\noindent\fbox{\parbox{0.98\linewidth}{%
\textbf{This study is:} a tokenizer-training and \emph{compression} benchmark; a
controlled comparison of corpus composition; a held-out technical-documentation
evaluation; a study of domain/general tradeoffs; a CPU-only reproducible artifact.

\smallskip
\textbf{This study is not:} an LLM quality evaluation; a model-adaptation method; a
serving or GPU benchmark; a billing study; an endorsement of replacing production
tokenizers; or a claim that fewer tokens always means a better system.%
}}
